\section{Human Validation}
\label{sec:human}

NLI-based faithfulness scoring captures entailment between answer claims
and retrieved context, but reviewers reasonably question whether the
metric tracks the kind of faithfulness humans actually care about. This
section reports a Prolific study designed to (a) calibrate NLI scores
against human judgement and (b) test whether the coherence paradox
survives when faithfulness is judged by humans rather than a model.

\subsection{Study design}

We sample 500 (case, condition) items stratified across the five
head-to-head conditions (§\ref{sec:headtohead}) — 100 items per
condition, balanced within each condition across the six datasets in
proportion to the multidataset query mix. Each item is double-annotated
by independent Prolific workers (no overlap), yielding 1{,}000 annotation
tasks. The full study budget is approximately £450 at the standard UK
Prolific rate (£9/hour, 1.5~min per item).

Annotators see the question, the retrieved passages, and the candidate
answer; they do \emph{not} see the condition label or any system
metadata. The rubric (full text in
\texttt{results/human\_validation/study\_v1/annotation\_guide.md}) asks
for:

\begin{itemize}
  \item A binary \textbf{Faithful} judgement (every claim supported by
        the passages, yes/no), with explicit handling of honest
        abstention (``I cannot find this in the passages'' = Faithful=YES
        when correct).
  \item A 1--5 \textbf{Correctness} rating.
  \item A 1--5 \textbf{Helpfulness} rating.
  \item An optional one-sentence comment.
\end{itemize}

Workers who fail a 3-item calibration set with $<$80\% agreement against
the gold rubric are rejected and the items re-issued.

\subsection{Inter-annotator agreement}

\begin{table}[!htbp]
\centering
\caption{Inter-annotator agreement (Fleiss' $\kappa$) on the 500-item
sample. Conventional bands: $<$0.20 poor, 0.21--0.40 fair, 0.41--0.60
moderate, 0.61--0.80 substantial, $>$0.81 almost perfect.}
\label{tab:kappa}
\small
\begin{tabular}{lc}
\toprule
Dimension                 & Fleiss' $\kappa$ \\
\midrule
Faithful (binary)         & TBD \\
Correctness (1--5 ord.)   & TBD \\
Helpfulness (1--5 ord.)   & TBD \\
\bottomrule
\end{tabular}
\end{table}

\subsection{NLI vs human}

\begin{table}[!htbp]
\centering
\caption{Correlation between NLI faithfulness score and human ratings
(Spearman / Pearson). $n =$ items with at least 2 annotations and a
defined NLI score.}
\label{tab:nli_human}
\small
\begin{tabular}{llcccc}
\toprule
NLI signal & Human dimension & $n$ & $\rho$ (Sp.) & $r$ (Pear.) & $p_{\rho}$ \\
\midrule
faithfulness\_score & human\_faithful\_rate & TBD & TBD & TBD & TBD \\
faithfulness\_score & mean\_correctness     & TBD & TBD & TBD & TBD \\
faithfulness\_score & mean\_helpfulness     & TBD & TBD & TBD & TBD \\
\bottomrule
\end{tabular}
\end{table}

\subsection{Does the paradox survive human evaluation?}

The load-bearing question is whether HCPC-v1 — which raises NLI-measured
faithfulness above baseline on parts of the §\ref{sec:results} table
while \emph{lowering} answer quality — also looks worse to humans, or
whether the paradox is an artifact of the NLI scorer.

\begin{table}[!htbp]
\centering
\caption{Per-condition human-faithfulness rate vs NLI-faithfulness rate.
The coherence paradox is preserved under human evaluation if
\texttt{hcpc\_v1} shows a lower \texttt{human\_faith\_rate} than
\texttt{baseline} on the same query set.}
\label{tab:paradox_human}
\small
\begin{tabular}{lcccc}
\toprule
Condition & $n$ & Human-faith rate & NLI-faith rate & Mean correctness \\
\midrule
baseline   & TBD & TBD & TBD & TBD \\
hcpc\_v1   & TBD & TBD & TBD & TBD \\
hcpc\_v2   & TBD & TBD & TBD & TBD \\
crag       & TBD & TBD & TBD & TBD \\
selfrag    & TBD & TBD & TBD & TBD \\
\bottomrule
\end{tabular}
\end{table}

\subsection{Reporting commitments}

We commit in advance to reporting the human-evaluation outcome regardless
of direction:

\begin{itemize}
  \item If $\rho_{\text{NLI},\text{human}} \geq 0.5$ \emph{and} the
        paradox replicates (HCPC-v1 human-faith $<$ baseline human-faith),
        the §\ref{sec:results} story stands.
  \item If $\rho_{\text{NLI},\text{human}} < 0.3$, we narrow the
        faithfulness claims to ``NLI-measured faithfulness'' throughout
        and treat the human study as evidence that NLI is the wrong
        oracle.
  \item If the paradox does \emph{not} replicate under humans (HCPC-v1
        human-faith $\geq$ baseline human-faith), we report it as a
        falsification: the §\ref{sec:results} effect was an NLI
        artifact and HCPC-v2's recovery would need a different
        justification.
\end{itemize}

The annotation pipeline, rubric, and analysis script are released as
part of ContextCoherenceBench (§\ref{sec:benchmark}) so the study can be
replicated or extended on additional conditions.
